\documentclass{situationist-preprint}

\title{Reliability Without Causal Contact}
\msid{MTH-02-b}
\msfamily{MTH-02}
\msdiscipline{Philosophy of Mathematics}
\mstype{Companion manuscript}
\mskeywords{Benacerraf's problem; reliabilism; safety; causal theory of knowledge; knowledge of the future}

\begin{document}
\maketitle

\begin{abstract}
MTH-02-a shows that derivation from true axioms is a reliable method and that reliable, justified belief is knowledge, so that abstract objects are known without causal contact. A residual worry is that reliability itself must be causally underwritten: a method is reliable only because the facts it tracks act on it. We argue that reliability is a modal property whose ground need not be causal, and that the worry, if sound, would also exclude knowledge of the future and of general laws, which the causal requirement was never meant to exclude.
\end{abstract}

\section{Reliability is modal}\label{s:modal}

By Definition~2.1 of MTH-02-a, a method is reliable if all its outputs in close cases are true; by Theorem~3.2, a justified belief formed by such a method is knowledge. Neither statement mentions causation. The worry is that this is an omission: in every ordinary case of reliable belief, a perceptual or instrumental one, the method is reliable because the fact believed causes the belief, and a reliable method with no causal link to its subject matter would be a coincidence too large to accept.

\begin{proposition}\label{p:ground}
A method can be reliable in virtue of a non-causal ground: a necessary connection between the conditions under which it issues an output and the truth of that output.
\end{proposition}

\begin{proof}
Let a method issue a belief only when the belief's content is derived by truth-preserving rules from true axioms. That the rules preserve truth is a necessary fact about them (Lemma~3.1 of MTH-02-a); that the axioms are true is fixed in every close case, since mathematical truths do not vary across close cases. So in every close case the method's outputs are true. What secures this is the necessary connection between derivability and truth, not any action of the facts on the believer.
\end{proof}

The ground is not a coincidence, because nothing is left to chance: the method does not happen to agree with the facts, it is constructed so that it cannot fail to. A thermometer is reliable because temperature acts on mercury; a correct derivation is reliable because of what derivability is.

\section{Two parallels}\label{s:par}

The causal requirement of \citet{Goldman1967} was designed for empirical knowledge, and \citet{Steiner1973} pointed out that it has difficulty even there. Two cases show that reliability without causal contact is not peculiar to mathematics.

\emph{The future.} An astronomer knows that a solar eclipse will occur on a given date. The eclipse does not cause her belief, since causes do not follow their effects. Her belief is formed by deriving the date from laws and present conditions, and the method is reliable because the laws hold in close cases. Goldman's theory accommodated such cases only by allowing the belief and the fact to have a common cause; there is no common cause of a belief and the infinitude of the primes, but there is no need of one, since the derivation does the work that the common cause does in the astronomical case.

\emph{General laws.} Knowledge that all copper conducts electricity concerns infinitely many cases, most of which have no causal relation to the knower. It is reliable because the inductive method is sensitive to the nomic structure the cases share, and that structure is not an object that acts on anyone.

\begin{objection}[Cheap reliability]
Every method that outputs necessary truths is trivially reliable, including guessing that happens to land on truths. Reliability without causal contact is therefore too easily had to explain knowledge \citep{ClarkeDoane2017}.
\end{objection}

\begin{reply}
Safety, as MTH-02-a uses it, concerns the method across the propositions it would be applied to in close cases, not the single proposition believed \citep{Williamson2000,Pritchard2005}. Guessing is unsafe because in close cases it outputs falsehoods about other propositions. Derivation is safe because no correct derivation in any close case outputs a falsehood. The condition is not cheap: it is met by proof and calculation, and failed by guessing and by careless calculation.
\end{reply}

\section{Conclusion}

Reliability is a modal property, and its ground can be a necessary connection rather than a causal one. That is how proof secures knowledge of abstract objects, and how prediction and induction secure knowledge of the future and of laws. Benacerraf's causal requirement excludes all three, and should be abandoned for all three.

\begin{thebibliography}{9}\small
\bibitem[Clarke-Doane(2017)]{ClarkeDoane2017} Clarke-Doane, J. (2017). What is the Benacerraf problem? In F.~Pataut (ed.), \emph{Truth, Objects, Infinity: New Perspectives on the Philosophy of Paul Benacerraf}, Springer.
\bibitem[Goldman(1967)]{Goldman1967} Goldman, A.~I. (1967). A causal theory of knowing. \emph{The Journal of Philosophy}, 64(12), 357--372.
\bibitem[Pritchard(2005)]{Pritchard2005} Pritchard, D. (2005). \emph{Epistemic Luck}. Oxford University Press.
\bibitem[Steiner(1973)]{Steiner1973} Steiner, M. (1973). Platonism and the causal theory of knowledge. \emph{The Journal of Philosophy}, 70(3), 57--66.
\bibitem[Williamson(2000)]{Williamson2000} Williamson, T. (2000). \emph{Knowledge and Its Limits}. Oxford University Press.
\end{thebibliography}

\end{document}
